\clearpage
\section{Daftar Repositori Kode Sumber dan Alias}

Berikut adalah daftar repositori kode sumber yang digunakan dalam penelitian ini. Setiap repositori mencakup dokumentasi lengkap dan panduan penggunaan. Kode panggilan (aliases) yang digunakan dalam naskah merujuk pada daftar ini untuk memudahkan penelusuran referensi teknis.

\begin{enumerate}
    \item \textbf{[Repo-Org-Skripsi]}: Organisasi GitHub yang menaungi seluruh repositori terkait penelitian tugas akhir ini.\\
    \url{https://github.com/Skripsi-Maulana-Anjari}
    
    \item \textbf{[Repo-Naskah]}: Repositori kode sumber LaTeX untuk penyusunan dokumen naskah skripsi lengkap.\\
    \url{https://github.com/Skripsi-Maulana-Anjari/naskah-makalah}
    
    \item \textbf{[Repo-Paper]}: Repositori kode sumber LaTeX untuk penyusunan makalah publikasi ilmiah.\\
    \url{https://github.com/Skripsi-Maulana-Anjari/naskah-latex}

    \item \textbf{[Repo-Gen-PoA]}: Modifikasi pada klien Geth dan alat bantu pembangkitan genesis block yang disesuaikan untuk konsensus Proof of Authority (PoA).\\
    \url{https://github.com/Skripsi-Maulana-Anjari/DChain-PoA-Geth}

    \item \textbf{[Repo-Deploy-PoA]}: Kumpulan skrip deployment dan manifest Kubernetes untuk implementasi jaringan blockchain dengan mekanisme konsensus Proof of Authority (PoA).\\
    \url{https://github.com/Skripsi-Maulana-Anjari/DChain-PoA-Deployer}

    \item \textbf{[Repo-Gen-PoS]}: Konfigurasi dan setup untuk klien Geth dan Prysm dalam implementasi konsensus Proof of Stake (PoS).\\
    \url{https://github.com/Skripsi-Maulana-Anjari/DChain-PoS-Geth}

    \item \textbf{[Repo-Benchmark-Pipeline]}: Ruang kerja (\textit{workspace}) untuk pipeline pengujian kinerja menggunakan Hyperledger Caliper beserta skrip definisi beban kerja (\textit{workload scripts}).\\
    \url{https://github.com/Skripsi-Maulana-Anjari/caliper-benchmark-workspace}

    \item \textbf{[Repo-Monitor]}: Konfigurasi dan skrip deployment untuk sistem pemantauan kinerja menggunakan Prometheus dan Grafana.\\
    \url{https://github.com/Skripsi-Maulana-Anjari/monitoring-tools}

    \item \textbf{[Repo-Benchmark-Results]}: Repositori penyimpanan data mentah hasil pengujian benchmark, mencakup log sistem, file CSV, dan laporan HTML.\\
    \url{https://github.com/Skripsi-Maulana-Anjari/benchmark-results}

    \item \textbf{[Repo-Analysis]}: Koleksi skrip Python yang digunakan untuk pemrosesan dan analisis data hasil pengujian serta visualisasi grafik.\\
    \url{https://github.com/Skripsi-Maulana-Anjari/analysis}
\end{enumerate}
